\chapter{\tr{Análise descritiva}{Descriptive analysis}}
\label{cap:descritiva}

\pt{%
Este capítulo descreve as métricas que os modelos dos capítulos seguintes vão usar: a exposição, o desfecho, a concentração dos acidentes entre veículos e unidades, o padrão no tempo, os tipos e causas de acidente, os eventos de telemetria e as fontes disponíveis para estimar a subnotificação. O objetivo é duplo: conhecer a distribuição de cada métrica antes de modelá-la e justificar, com os dados, as escolhas do Capítulo~\ref{cap:metodos}. A análise preditiva (Capítulo~\ref{cap:preditiva}) e a prescritiva (Capítulo~\ref{cap:prescritiva}) partem daqui.
}
\en{%
This chapter describes the metrics that the models in the following chapters will use: exposure, the outcome, the concentration of accidents across vehicles and sites, the pattern over time, accident types and causes, telematics events and the sources available to estimate underreporting. The aim is twofold: to know the distribution of each metric before modeling it and to justify, with the data, the choices in Chapter~\ref{cap:metodos}. The predictive (Chapter~\ref{cap:preditiva}) and prescriptive (Chapter~\ref{cap:prescritiva}) analyses build on it.
}

\section{\tr{Convenções}{Conventions}}

\pt{%
Os dados são os da operação descritos no Capítulo~\ref{cap:dados}, extraídos só como agregados por veículo e mês ou por categoria, sem dados pessoais. O painel cobre doze meses comuns à telemetria e aos relatos e inclui os veículo-meses com ao menos dez horas de viagem; os relatos e as fontes auxiliares cobrem o período disponível de cada uma. Para preservar a confidencialidade, todas as figuras e números do capítulo são relativos: proporções, índices em relação à média, razões e exposição em relação à mediana. As unidades não são identificadas. As figuras são geradas por \texttt{analise/real/descritiva.py} (Python) e \texttt{r/descritiva.R} (R), e a Tabela~\ref{tab:descritiva} resume os números citados.
}
\en{%
The data are the operational data described in Chapter~\ref{cap:dados}, extracted only as aggregates per vehicle and month or per category, with no personal data. The panel covers twelve months common to telematics and reports and includes vehicle-months with at least ten trip hours; reports and auxiliary sources cover each one's available period. To preserve confidentiality, all figures and numbers in the chapter are relative: proportions, indices relative to the mean, ratios and exposure relative to the median. Sites are not identified. Figures are generated by \texttt{analise/real/descritiva.py} (Python) and \texttt{r/descritiva.R} (R), and Table~\ref{tab:descritiva} summarizes the numbers cited.
}

\section{\tr{Exposição}{Exposure}}
\label{sec:desc-exposicao}

\begin{figure}[!htbp]
\centering
\includegraphics[width=\textwidth]{\arqfig{fig_desc_exposicao}}
\caption{\tr%
  {Exposição. (a) a (c): distribuição por veículo-mês, por tipo de uso, relativa à mediana geral (escala log; cada ponto é um veículo-mês). (d) Fração das horas pagas que são horas extras, por unidade e mês da folha; linha tracejada: mediana.}%
  {Exposure. (a) to (c): distribution per vehicle-month, by use type, relative to the overall median (log scale; each point is a vehicle-month). (d) Share of paid hours that are overtime, per payroll site and month; dashed line: median.}}
\label{fig:desc-exposicao}
\end{figure}

\pt{%
As horas de viagem por veículo-mês têm distribuição parecida nos três tipos de uso, com mediana próxima da geral e dispersão grande dentro de cada tipo (Figura~\ref{fig:desc-exposicao}a). As milhas, ao contrário, separam os tipos por uma ordem de grandeza (Figura~\ref{fig:desc-exposicao}b): o spotter anda a cerca de \descVelSpotter{} da velocidade média do veículo local, e o shuttle, a cerca de \descVelShuttle{} (Figura~\ref{fig:desc-exposicao}c). É a forma visual do argumento do Capítulo~\ref{cap:metodos}: usar a distância como exposição faria o spotter parecer muito mais arriscado só porque quase não roda, e por isso a exposição principal é o tempo. Na folha de pagamento, as horas extras representam em mediana cerca de \descHorasExtras\% das horas de cada unidade-mês, com muitas unidades-mês sem nenhuma e uma cauda longa (Figura~\ref{fig:desc-exposicao}d). Essa variação é o que torna a fração de horas extras uma covariável candidata \parencite{dembe2005}.
}
\en{%
Trip hours per vehicle-month have a similar distribution across the three use types, with a median close to the overall one and a large spread within each type (Figure~\ref{fig:desc-exposicao}a). Miles, in contrast, separate the types by an order of magnitude (Figure~\ref{fig:desc-exposicao}b): the spotter moves at about \descVelSpotter{} of the local vehicle's average speed, and the shuttle at about \descVelShuttle{} (Figure~\ref{fig:desc-exposicao}c). This is the visual form of the argument in Chapter~\ref{cap:metodos}: using distance as exposure would make the spotter look much riskier only because it barely moves, which is why the main exposure is time. In payroll, overtime accounts for a median of about \descHorasExtras\% of each site-month's hours, with many site-months with none and a long tail (Figure~\ref{fig:desc-exposicao}d). This variation is what makes the overtime share a candidate covariate \parencite{dembe2005}.
}

\section{\tr{Desfecho: contagens e superdispersão}{Outcome: counts and overdispersion}}
\label{sec:desc-desfecho}

\begin{figure}[!htbp]
\centering
\includegraphics[width=\textwidth]{\arqfig{fig_desc_contagem}}
\caption{\tr%
  {Desfecho. (a) Proporção de veículo-meses com 0, 1, 2 e 3 ou mais acidentes (barras) e o esperado por três modelos com tipo de uso e horas de viagem como \textit{offset} (pontos; média de 300 simulações). (b) O mesmo para o total de acidentes de cada veículo no período. (c) Risco relativo estimado de cada veículo (efeito aleatório da binomial negativa), com intervalo de 95\%. Escala raiz em (a) e (b), log em (c).}%
  {Outcome. (a) Share of vehicle-months with 0, 1, 2 and 3 or more accidents (bars) and the expected share under three models with use type and trip hours as offset (points; mean of 300 simulations). (b) The same for each vehicle's total accidents in the period. (c) Estimated relative risk of each vehicle (negative binomial random effect), with 95\% interval. Square-root scale in (a) and (b), log in (c).}}
\label{fig:desc-contagem}
\end{figure}

\pt{%
O acidente é raro no nível do veículo-mês: cerca de \descZeros\% dos veículo-meses não têm nenhum, e a variância das contagens é apenas \descVarMedia{} vez a média. Nesse nível, o próprio Poisson com exposição e tipo de uso já reproduz a proporção de zeros (\descRzerosPoisson\% esperados contra \descRzerosObs\% observados; Figura~\ref{fig:desc-contagem}a), o que confirma que não há excesso de zeros a modelar \parencite{lord2005}. A heterogeneidade aparece quando se soma o período de cada veículo (Figura~\ref{fig:desc-contagem}b): o Poisson espera \descRveicZeroPoisson\% de veículos sem acidente, contra \descRveicZeroObs\% observados, e a binomial negativa com efeito de veículo acerta essa proporção (\descRveicZeroRE\%). O desvio padrão do efeito de veículo é de \descRdpVeiculo{} na escala log, mas os intervalos de cada veículo são largos e quase todos cruzam 1 (Figura~\ref{fig:desc-contagem}c): com um ano de dados, dá para afirmar que os veículos diferem, mas não para dizer com segurança quais são os piores. É o cenário em que o encolhimento é necessário.
}
\en{%
Accidents are rare at the vehicle-month level: about \descZeros\% of vehicle-months have none, and the variance of the counts is only \descVarMedia{} times the mean. At this level, the Poisson model with exposure and use type already reproduces the share of zeros (\descRzerosPoisson\% expected against \descRzerosObs\% observed; Figure~\ref{fig:desc-contagem}a), which confirms there is no excess of zeros to model \parencite{lord2005}. Heterogeneity appears when each vehicle's period is summed (Figure~\ref{fig:desc-contagem}b): the Poisson model expects \descRveicZeroPoisson\% of vehicles with no accident, against \descRveicZeroObs\% observed, and the negative binomial with a vehicle effect matches this share (\descRveicZeroRE\%). The standard deviation of the vehicle effect is \descRdpVeiculo{} on the log scale, but each vehicle's interval is wide and almost all cross 1 (Figure~\ref{fig:desc-contagem}c): with one year of data, one can claim that vehicles differ, but not say with confidence which ones are the worst. This is the setting in which shrinkage is needed.
}

\section{\tr{Concentração e diferenças entre unidades}{Concentration and differences across sites}}
\label{sec:desc-concentracao}

\begin{figure}[!htbp]
\centering
\includegraphics[width=\textwidth]{\arqfig{fig_desc_concentracao}}
\caption{\tr%
  {Concentração. (a) e (b) Curvas de concentração: fração acumulada dos acidentes (linha cheia) e das horas (tracejada) contra a fração acumulada de veículos ou unidades, ordenados do maior para o menor; faixa: 95\% das curvas simuladas sob Poisson com taxa constante dentro de cada tipo de uso. (c) Razão entre acidentes observados e esperados por unidade (mesmo modelo), com intervalo exato de 95\%, para as unidades com ao menos um acidente esperado; \corpb{em vermelho}{em preto}, as que excluem 1. (d) Horas de viagem por unidade relativas à mediana.}%
  {Concentration. (a) and (b) Concentration curves: cumulative share of accidents (solid line) and hours (dashed) against the cumulative share of vehicles or sites, ordered from largest to smallest; band: 95\% of curves simulated under Poisson with a constant rate within each use type. (c) Ratio of observed to expected accidents per site (same model), with exact 95\% interval, for sites with at least one expected accident; \corpb{in red}{in black}, those that exclude 1. (d) Trip hours per site relative to the median.}}
\label{fig:desc-concentracao}
\end{figure}

\pt{%
Os acidentes parecem muito concentrados: metade dos veículos não teve nenhum, e o coeficiente de Gini dos acidentes entre veículos é de \descGiniObs{} (Figura~\ref{fig:desc-concentracao}a). Mas quase toda essa concentração é a esperada por acaso e exposição: um Poisson com taxa constante dentro de cada tipo de uso produz Gini de \descGiniPoisson{}, e a curva observada fica só um pouco acima da faixa simulada. Entre unidades, o padrão se repete (Gini de \descGiniSiteObs{} contra \descGiniSitePoisson{}; Figura~\ref{fig:desc-concentracao}b). A consequência prática é importante para a gestão: ver que poucos veículos ou unidades concentram os acidentes não é, por si só, evidência de que eles sejam mais arriscados.
}
\en{%
Accidents look highly concentrated: half of the vehicles had none, and the Gini coefficient of accidents across vehicles is \descGiniObs{} (Figure~\ref{fig:desc-concentracao}a). But almost all of this concentration is what chance and exposure would produce: a Poisson model with a constant rate within each use type yields a Gini of \descGiniPoisson{}, and the observed curve lies only slightly above the simulated band. Across sites the pattern repeats (Gini of \descGiniSiteObs{} against \descGiniSitePoisson{}; Figure~\ref{fig:desc-concentracao}b). The practical consequence matters for management: seeing that a few vehicles or sites concentrate accidents is not, in itself, evidence that they are riskier.
}

\pt{%
As unidades diferem muito em tamanho: a que está no percentil 90 da exposição opera cerca de \descPnoventaPdez{} vezes mais horas que a do percentil 10 (Figura~\ref{fig:desc-concentracao}d). Por isso a contagem bruta mede sobretudo o tamanho. Ajustando pela exposição e pelo tipo de uso, cerca de \descForaIC\% das unidades com ao menos um acidente esperado ficam com intervalo de 95\% inteiramente acima ou abaixo de 1 (Figura~\ref{fig:desc-concentracao}c). A proporção é menor que a da exploração do Capítulo~\ref{cap:viabilidade}, que não ajustava pelo tipo de uso: parte da diferença aparente entre unidades vem da mistura de operações de cada uma. O que sobra é a variação que o efeito de unidade do modelo hierárquico deve capturar (H1b).
}
\en{%
Sites differ greatly in size: the one at the 90th percentile of exposure operates about \descPnoventaPdez{} times more hours than the one at the 10th percentile (Figure~\ref{fig:desc-concentracao}d). That is why the raw count mostly measures size. Adjusting for exposure and use type, about \descForaIC\% of sites with at least one expected accident have a 95\% interval entirely above or below 1 (Figure~\ref{fig:desc-concentracao}c). The share is smaller than in the exploration of Chapter~\ref{cap:viabilidade}, which did not adjust for use type: part of the apparent difference across sites comes from each site's mix of operations. What remains is the variation the hierarchical model's site effect should capture (H1b).
}

\section{\tr{Padrão no tempo}{Pattern over time}}
\label{sec:desc-tempo}

\begin{figure}[!htbp]
\centering
\includegraphics[width=\textwidth]{\arqfig{fig_desc_tempo}}
\caption{\tr%
  {Tempo. (a) Relatos de acidente por mês relativos à média mensal; hachurado: relatos pendentes de revisão no período de implantação do aplicativo, suspeitos de teste. (b) e (c) Distribuição dos acidentes (sem os suspeitos) por hora do dia e por dia da semana; linha pontilhada: distribuição uniforme; faixas: turnos de 8 horas. (d) Horas da folha por semana e horas de viagem por mês, relativas à média de cada série.}%
  {Time. (a) Accident reports per month relative to the monthly mean; hatched: reports pending review during the application roll-out, suspected tests. (b) and (c) Distribution of accidents (excluding suspected tests) by hour of day and day of week; dotted line: uniform distribution; bands: 8-hour shifts. (d) Payroll hours per week and trip hours per month, relative to each series' mean.}}
\label{fig:desc-tempo}
\end{figure}

\pt{%
O número de relatos varia bastante de um mês para outro (Figura~\ref{fig:desc-tempo}a), e cerca de \descSuspeitos\% dos relatos de acidente são suspeitos de teste: estavam pendentes de revisão e foram enviados durante a implantação do aplicativo, concentrados entre maio e outubro de 2026. A regra de exclusão desses relatos muda o desfecho de forma desigual no tempo e precisa de análise de sensibilidade, como previsto no Capítulo~\ref{cap:metodos}.
}
\en{%
The number of reports varies considerably from one month to the next (Figure~\ref{fig:desc-tempo}a), and about \descSuspeitos\% of accident reports are suspected tests: they were pending review and submitted during the application roll-out, concentrated between May and October 2026. The rule for excluding these reports changes the outcome unevenly over time and needs a sensitivity analysis, as planned in Chapter~\ref{cap:metodos}.
}

\pt{%
Quase todos os relatos (\descComHora\%) trazem a hora do acidente. Os acidentes se distribuem pelas 24 horas, com \descTurnoUm\%, \descTurnoDois\% e \descTurnoTres\% nos turnos da manhã, da tarde e da noite, e um pico no início da noite (Figura~\ref{fig:desc-tempo}b). Por dia da semana, há mais acidentes de terça a quinta e menos no fim de semana (Figura~\ref{fig:desc-tempo}c), o que acompanha o volume de operação. Sem a exposição por hora e por dia, essas distribuições ainda não podem ser lidas como risco; servem para mostrar que um modelo por turno é viável. Já a exposição cresceu de forma contínua ao longo do ano, tanto nas horas da folha quanto nas horas de viagem (Figura~\ref{fig:desc-tempo}d), com a entrada de novas unidades e veículos: comparar meses sem \textit{offset} e sem efeito de calendário confundiria crescimento da operação com aumento de risco.
}
\en{%
Almost all reports (\descComHora\%) include the time of the accident. Accidents are spread over the 24 hours, with \descTurnoUm\%, \descTurnoDois\% and \descTurnoTres\% in the morning, afternoon and night shifts, and a peak in the early evening (Figure~\ref{fig:desc-tempo}b). By day of week, there are more accidents from Tuesday to Thursday and fewer on weekends (Figure~\ref{fig:desc-tempo}c), which follows operating volume. Without exposure by hour and by day, these distributions cannot yet be read as risk; they show that a shift-level model is feasible. Exposure, in turn, grew steadily over the year, in both payroll hours and trip hours (Figure~\ref{fig:desc-tempo}d), as new sites and vehicles came in: comparing months without an offset and a calendar effect would confuse operational growth with rising risk.
}

\section{\tr{Tipos e causas de acidente}{Accident types and causes}}
\label{sec:desc-tipos}

\begin{figure}[!htbp]
\centering
\includegraphics[width=\textwidth]{\arqfig{fig_desc_tipos}}
\caption{\tr%
  {Tipos e causas. (a) Distribuição dos relatos por tipo. (b) Participação de cada tipo de uso nos acidentes, nas horas de viagem e nas milhas. (c) Causa atribuída aos acidentes, em categorias fechadas do formulário de relato.}%
  {Types and causes. (a) Distribution of reports by type. (b) Share of each use type in accidents, trip hours and miles. (c) Cause assigned to accidents, in closed categories of the report form.}}
\label{fig:desc-tipos}
\end{figure}

\pt{%
Os acidentes com veículo são cerca de \descAcidentesPct\% dos relatos; lesões e outros eventos são minoria (Figura~\ref{fig:desc-tipos}a), o que reforça a escolha de todas as colisões como desfecho principal. O spotter responde pela maior parte dos acidentes, por pouco mais da metade das horas e por uma fração pequena das milhas; o veículo local, pelo inverso (Figura~\ref{fig:desc-tipos}b). As causas mais frequentes são as do acidente de baixa velocidade típico do pátio: problemas na porta da carreta e na fixação de equipamento, erro operacional e ré com folga mal avaliada somam cerca de \descTopCausas\% dos acidentes, seguidas de colisão com objeto fixo, carreta caída e problemas de engate (Figura~\ref{fig:desc-tipos}c). Colisões com outros veículos, o acidente típico da estrada, são minoria. É um retrato que antecipa por que eventos de telemetria desenhados para a estrada podem carregar pouco sinal no pátio (H2).
}
\en{%
Vehicle accidents are about \descAcidentesPct\% of reports; injuries and other events are a minority (Figure~\ref{fig:desc-tipos}a), which supports choosing all collisions as the main outcome. The spotter accounts for most accidents, a little over half of the hours and a small share of the miles; the local vehicle, for the opposite (Figure~\ref{fig:desc-tipos}b). The most frequent causes are those of the low-speed accident typical of the yard: trailer door and equipment securing problems, operational error and backing with misjudged clearance add up to about \descTopCausas\% of accidents, followed by striking a fixed object, dropped trailers and coupling problems (Figure~\ref{fig:desc-tipos}c). Collisions with other vehicles, the typical road accident, are a minority. This picture anticipates why telematics events designed for the road may carry little signal in the yard (H2).
}

\section{\tr{Eventos de telemetria}{Telematics events}}
\label{sec:desc-eventos}

\begin{figure}[!htbp]
\centering
\includegraphics[width=\textwidth]{\arqfig{fig_desc_eventos}}
\caption{\tr%
  {Eventos de telemetria. (a) Proporção de veículo-meses com ao menos um evento de cada tipo. (b) Correlação de Spearman entre as taxas de eventos por hora de viagem, a taxa de acidentes por hora e as horas de viagem, por veículo-mês.}%
  {Telematics events. (a) Share of vehicle-months with at least one event of each type. (b) Spearman correlation between event rates per trip hour, the accident rate per hour and trip hours, per vehicle-month.}}
\label{fig:desc-eventos}
\end{figure}

\pt{%
A maioria dos tipos de evento é rara: em boa parte dos veículo-meses não há nenhum, e só a ausência de cinto aparece na maioria deles (Figura~\ref{fig:desc-eventos}a). As taxas dos diferentes tipos são pouco correlacionadas entre si; a maior correlação, de \descRcorMax{}, é entre frenagem brusca e curva brusca, dois eventos cinemáticos. A correlação de cada tipo com a taxa de acidentes do mesmo veículo-mês não passa de \descRcorAcidMax{} em módulo, e com as horas de viagem, de \descRcorHorasMax{} (Figura~\ref{fig:desc-eventos}b). Duas consequências para o Capítulo~\ref{cap:preditiva}: testar cada tipo num modelo próprio não perde muito por ignorar os demais, e o sinal marginal da telemetria é fraco já antes de qualquer ajuste, o que reforça a necessidade dos testes de equivalência de H2b.
}
\en{%
Most event types are rare: many vehicle-months have none, and only the no-seatbelt event appears in most of them (Figure~\ref{fig:desc-eventos}a). The rates of different types are weakly correlated with each other; the largest correlation, \descRcorMax{}, is between harsh braking and harsh turns, two kinematic events. Each type's correlation with the accident rate of the same vehicle-month does not exceed \descRcorAcidMax{} in absolute value, and with trip hours, \descRcorHorasMax{} (Figure~\ref{fig:desc-eventos}b). Two consequences for Chapter~\ref{cap:preditiva}: testing each type in its own model loses little by ignoring the others, and the marginal telematics signal is weak even before any adjustment, which reinforces the need for the equivalence tests of H2b.
}

\section{\tr{Fontes para a subnotificação}{Sources for underreporting}}
\label{sec:desc-fontes}

\begin{figure}[!htbp]
\centering
\includegraphics[width=\textwidth]{\arqfig{fig_desc_fontes}}
\caption{\tr%
  {Fontes. (a) Relatos de acidente, ordens de manutenção por dano e colisões detectadas pela telemetria, por mês, relativos à média de cada série no período comum. (b) Ordens de manutenção por dano relativas à média mensal; tons: grafia antiga e nova do motivo do reparo; linha tracejada: mudança de regime. (c) Nota das auditorias de segurança por unidade (\% dos pontos possíveis); linha tracejada: mediana. (d) Notas por mês.}%
  {Sources. (a) Accident reports, damage repair orders and telematics-detected collisions per month, relative to each series' mean over the common period. (b) Damage repair orders relative to the monthly mean; shades: old and new coding of the repair reason; dashed line: regime change. (c) Site safety audit scores (\% of possible points); dashed line: median. (d) Scores by month.}}
\label{fig:desc-fontes}
\end{figure}

\pt{%
As três fontes de acidentes andam de forma diferente ao longo do ano (Figura~\ref{fig:desc-fontes}a). As colisões detectadas pela telemetria crescem quase continuamente, acompanhando a frota, enquanto os relatos e as ordens de dano oscilam sem seguir umas às outras mês a mês, o que já sugere que as fontes captam universos diferentes. As ordens de manutenção por dano mudam de regime em novembro de 2025 (Figura~\ref{fig:desc-fontes}b): a grafia do motivo muda, o volume passa a ser de duas a três vezes o anterior, e o motivo que indicava acidente já relatado praticamente desaparece. Para a Etapa~4, isso tem dois lados. No período comum com os relatos, a dúvida sobre esse motivo quase não pesa; em compensação, é preciso entender o que passou a ser registrado como dano antes de usar as ordens como segunda fonte \preencher{verificar com a área de manutenção o que mudou em novembro de 2025}.
}
\en{%
The three accident sources move differently over the year (Figure~\ref{fig:desc-fontes}a). Telematics-detected collisions grow almost continuously, following the fleet, while reports and damage orders fluctuate without tracking each other month by month, which already suggests that the sources capture different universes. Damage repair orders change regime in November 2025 (Figure~\ref{fig:desc-fontes}b): the coding of the reason changes, volume becomes two to three times the previous level, and the reason that flagged an already reported accident practically disappears. For Stage~4 this cuts both ways. In the period shared with reports, the doubt about that reason barely matters; on the other hand, one must understand what came to be recorded as damage before using the orders as a second source \preencher{check with the maintenance team what changed in November 2025}.
}

\pt{%
As auditorias de segurança por unidade, registradas desde meados de 2026, têm nota mediana de \descNotaAud\% dos pontos possíveis, com a maioria entre 90\% e 100\% (Figura~\ref{fig:desc-fontes}c e d). A nota é uma medida candidata do clima de segurança de cada unidade, que a literatura associa à propensão a relatar \parencite{probst2008,probst2010}, mas o efeito teto limita a variação disponível para explicar diferenças entre unidades.
}
\en{%
Site safety audits, recorded since mid-2026, have a median score of \descNotaAud\% of possible points, with most between 90\% and 100\% (Figure~\ref{fig:desc-fontes}c and d). The score is a candidate measure of each site's safety climate, which the literature associates with the propensity to report \parencite{probst2008,probst2010}, but the ceiling effect limits the variation available to explain differences across sites.
}

\section{\tr{Síntese}{Summary}}

\begin{table}[!htbp]
\centering\small
\caption{\tr{Resumo descritivo (só valores relativos).}{Descriptive summary (relative values only).}}
\label{tab:descritiva}
\begin{tabularx}{\textwidth}{X r}
\toprule
\tr{Métrica}{Metric} & \tr{Valor}{Value} \\
\midrule
\tabinput{tabelas/tab_descritiva.tex}
\bottomrule
\end{tabularx}
\end{table}

\pt{%
\begin{enumerate}
  \item \textbf{Exposição pelo tempo, ajuste pelo tipo de uso.} As milhas separam os tipos de uso por uma ordem de grandeza; as horas, não.
  \item \textbf{Binomial negativa com efeito de veículo, sem inflação de zeros.} O excesso de zeros não existe no veículo-mês; a heterogeneidade aparece entre veículos.
  \item \textbf{Concentração não é risco.} Quase toda a concentração dos acidentes em poucos veículos e unidades é a esperada por acaso e exposição; comparar exige taxas ajustadas e intervalos.
  \item \textbf{Calendário e limpeza.} A operação cresceu ao longo do ano, e um quarto dos relatos de acidente precisa de regra de exclusão documentada.
  \item \textbf{Turno é viável.} A hora do acidente está quase sempre preenchida, o que permite modelos por turno.
  \item \textbf{Telemetria com sinal marginal fraco.} Eventos raros, pouco correlacionados entre si e com os acidentes.
  \item \textbf{Fontes de subnotificação com ressalvas.} A manutenção mudou de regime, e as auditorias têm efeito teto.
\end{enumerate}
}
\en{%
\begin{enumerate}
  \item \textbf{Time-based exposure, adjustment for use type.} Miles separate use types by an order of magnitude; hours do not.
  \item \textbf{Negative binomial with a vehicle effect, no zero inflation.} There is no excess of zeros at the vehicle-month level; heterogeneity appears across vehicles.
  \item \textbf{Concentration is not risk.} Almost all the concentration of accidents in a few vehicles and sites is what chance and exposure produce; comparison requires adjusted rates and intervals.
  \item \textbf{Calendar and cleaning.} Operations grew over the year, and a quarter of accident reports needs a documented exclusion rule.
  \item \textbf{Shifts are feasible.} The accident time is almost always filled in, which allows shift-level models.
  \item \textbf{Weak marginal telematics signal.} Rare events, weakly correlated with each other and with accidents.
  \item \textbf{Underreporting sources with caveats.} Maintenance changed regime, and audits have a ceiling effect.
\end{enumerate}
}
